    \subsection{応用難易度}\label{節：応用難易度}\label{節：高難易度のもの}
      ここでは,標準問題で学んだ解法を複数組み合わせる10題を扱う.
      見た目の複雑さをそのまま計算するのではなく,共通の分母,和と差,積や総和の構造を探そう.
      どの数量に変換するかと,変換をどの順番で使うかを考えてほしい.
      以下,\,$n$は正の整数とし,必要なときは$S_n=\sum_{k=1}^{n}a_k$とおく.

      \subsubsection{問題}\label{節：応用難易度：問題}\label{節：高難易度のもの：問題}
        \begin{enumerate}[label=\textbf{問\arabic*.},leftmargin=3.8em,format=\bookitemlink{practice-advanced}{problem}{answer}]
          % A1: 階乗で重みをそろえた総和
          \item $\displaystyle a_1=1$ とし,
                \[
                  a_{n+1}
                  =4(n+1)a_n
                    -(n+1)!\sum_{k=1}^{n}\frac{a_k}{k!}
                \]
                で数列$\{a_n\}$を定める.一般項を求めよ.
          % A2: 添字を含む分母と不動点
          \item $\displaystyle a_1=2$ とし,
                \[
                  a_{n+1}
                   =\frac{(n+1)a_n+1-n}{na_n+2-n}
                \]
                で数列$\{a_n\}$を定める.一般項を求めよ.
          % A3: 三角形状の重みをもつ積
          \item $\displaystyle a_1=1$ とし,
                \[
                  a_{n+1}
                   =2^n\prod_{k=1}^{n}a_k^{\,n-k+1}
                \]
                で正の数列$\{a_n\}$を定める.一般項を求めよ.
          % A4: 非線形3項間式から階比を作る
          \item $\displaystyle a_1=1,\qquad a_2=2$ とし,
                \[
                  (na_n+a_{n+1})a_{n+2}
                    =(n+1)a_{n+1}^{\,2}
                \]
                で正の数列$\{a_n\}$を定める.一般項を,階乗を用いて表せ.
          % A5: 和と差で異なる構造が現れる連立
          \item $\displaystyle a_1=\frac32,\qquad b_1=\frac12$ とし,
                \begin{align*}
                 a_{n+1}
                  &=(n+1)(a_n+b_n)\\
                  &\quad+\frac{3(a_n-b_n)}
                      {2\{3+(n+1)(a_n-b_n)\}},\\
                 b_{n+1}
                  &=(n+1)(a_n+b_n)\\
                  &\quad-\frac{3(a_n-b_n)}
                      {2\{3+(n+1)(a_n-b_n)\}}
                \end{align*}
                で2つの数列を定める.それぞれの一般項を求めよ.
          % A6: 積の3項間関係と二重の階差
          \item $\displaystyle a_1=2,\qquad a_2=16$ とし,
                \[
                  a_{n+2}a_n=2^{2n+1}a_{n+1}^{\,2}
                \]
                を満たす正の数列$\{a_n\}$の一般項を求めよ.
          % A7: 添字とともに動く基準点
          \item $\displaystyle a_1=3$ とし,
                \[
                  a_{n+1}
                   =\frac{(n+1)^2a_n+n(n+1)}
                          {a_n+n(n+1)}
                \]
                で数列$\{a_n\}$を定める.一般項を求めよ.
          % A8: 係数の和が零になる3項間関係
          \item $\displaystyle a_1=0,\qquad a_2=1$ とし,
                \begin{align*}
                 &(n+2)(n+4)a_{n+2}\\
                 &\quad-(2n^2+10n+11)a_{n+1}\\
                 &\quad+(n+1)(n+3)a_n=0
                \end{align*}
                を満たす数列$\{a_n\}$の一般項を求めよ.
          % A9: 掛け算で結び付いた二つの数列
          \item $\displaystyle a_1=b_1=2$ とし,
                \[
                 \begin{cases}
                  a_{n+1}=2^{2n+1}a_n^{\,2}b_n,\\
                  b_{n+1}=\dfrac{a_nb_n^{\,2}}{2^{2n+1}}
                 \end{cases}
                \]
                で2つの正の数列を定める.それぞれの一般項を求めよ.
          % A10: 二つの平方が作る分数式
          \item $\displaystyle a_1=\frac97$ とし,
                \[
                  a_{n+1}
                   =\frac{(a_n+1)^2+2^n(a_n-1)^2}
                          {(a_n+1)^2-2^n(a_n-1)^2}
                \]
                で数列$\{a_n\}$を定める.一般項を求めよ.
        \end{enumerate}

      \subsubsection{方針}\label{節：応用難易度：方針}\label{節：高難易度のもの：方針}
        \begin{enumerate}[label=\textbf{問\arabic*.},leftmargin=3.8em,format=\bookitemlink{practice-advanced}{hint}{problem}]
          % A1: 階乗で重みをそろえた総和
          \item 総和の中の$a_k/k!$と,前項の係数$n+1$を同時に見よう.まず$b_n=a_n/n!$とおき,その部分和を消去する.得られる3項間漸化式では,さらに$2^{n-1}$で割る.
          % A2: 添字を含む分母と不動点
          \item 係数に$n$があるので,初めから一定の公比を探してもうまくいかない.まず$a_{n+1}-1$を計算し,その逆数を取る.最後に$n$の一次式を加えて等比型にする.
          % A3: 三角形状の重みをもつ積
          \item 対数を取ると,係数が$n-k+1$の重み付き総和になる.隣接する式の差を取ると重みが$1$になり,もう一度差を取ると総和が消える.
          % A4: 非線形3項間式から階比を作る
          \item 項そのものではなく,\, $r_n=a_{n+1}/a_n$を導入する.次に$1/r_n$の漸化式を作り,定数$1$を引くと積の相殺が見える.
          % A5: 和と差で異なる構造が現れる連立
          \item 2式に共通する部分と,符号だけが違う部分を分ける.和は係数が$n$に依存する階比型で,差は逆数を取ると階差型になる.
          % A6: 積の3項間関係と二重の階差
          \item 対数を取った後,\, $x_{n+1}-x_n$を新しい数列として扱う.また,対数を取る前に$a_{n+1}/a_n$を作るルートもある.
          % A7: 添字とともに動く基準点
          \item $a_n=n$を代入すると次項は$n+1$になる.一定の不動点を探す前に,\, $c_n=a_n/n$とおいて基準をそろえよう.その後$c_n=1,-1$からの距離の比を作る.
          % A8: 係数の和が零になる3項間関係
          \item 3つの係数を符号も含めて足すと$0$になる.式を隣接する項の差でまとめ,その階差数列を階比型として解く.最後の和には部分分数分解を使う.
          % A9: 掛け算で結び付いた二つの数列
          \item 対数を取ると,対称な連立漸化式になる.和と差を取れば,指数の増え方を分離できる.対数を取る前に$a_nb_n$と$a_n/b_n$を作る方法も考えよう.
          % A10: 二つの平方が作る分数式
          \item 分子と分母の和・差を使って$a_{n+1}\pm1$を計算する.それらの比を取ると平方が現れるが,係数$2^n$も残る.その次に対数を使う.
        \end{enumerate}

      \subsubsection{解答}\label{節：応用難易度：解答}\label{節：高難易度のもの：解答}
        \begin{enumerate}[label=\textbf{問\arabic*.},leftmargin=3.8em,format=\bookitemlink{practice-advanced}{answer}{problem}]
          % A1: 階乗で重みをそろえた総和
          \item 一般項は
                \[
                  a_n=n!(n+1)2^{n-2}
                \]
                である.この式は$n=1$でも成り立つ.
          % A2: 添字を含む分母と不動点
          \item 一般項は
                \[
                  a_n
                  =1+\frac{1}{3\cdot2^{n-1}-n-1}
                \]
                である.
          % A3: 三角形状の重みをもつ積
          \item $\displaystyle
                 \alpha=\frac{3+\sqrt5}{2},\quad
                 \beta=\frac{3-\sqrt5}{2}$
                とおくと,一般項は
                \[
                  a_n=2^{\frac{\alpha^{n-1}-\beta^{n-1}}{\sqrt5}}
                \]
                である.
          % A4: 非線形3項間式から階比を作る
          \item 一般項は
                \[
                  a_n
                   =\frac{4^{n-1}\{(n-1)!\}^2}{(2n-2)!}
                \]
                である.ここでは$0!=1$とする.
          % A5: 和と差で異なる構造が現れる連立
          \item 一般項は
                \begin{align*}
                 a_n&=2^{n-1}n!+\frac3{n^2+n+4},\\
                 b_n&=2^{n-1}n!-\frac3{n^2+n+4}
                \end{align*}
                である.
          % A6: 積の3項間関係と二重の階差
          \item $\displaystyle
                 E_n=\frac{n(n-1)(2n-1)}6+2n-1$
                とおくと,
                \[
                  a_n=2^{E_n}
                \]
                である.
          % A7: 添字とともに動く基準点
          \item 一般項は
                \[
                  a_n
                   =n\,\frac{n^2+n+1}{n^2+n-1}
                \]
                である.
          % A8: 係数の和が零になる3項間関係
          \item 一般項は
                \[
                  a_n
                   =\frac{10}{3}-\frac4{n+1}-\frac4{n+2}
                \]
                である.
          % A9: 掛け算で結び付いた二つの数列
          \item 一般項は
                \begin{align*}
                 a_n&=2^{\,3^{n-1}+n^2-1},\\
                 b_n&=2^{\,3^{n-1}-n^2+1}
                \end{align*}
                である.
          % A10: 二つの平方が作る分数式
          \item $\displaystyle E_n=n+1+2^{n-1}$ とおくと,一般項は
                \[
                  a_n=\frac{2^{E_n}+1}{2^{E_n}-1}
                \]
                である.
        \end{enumerate}

      \subsubsection{解法}\label{節：応用難易度：解法}\label{節：高難易度のもの：解法}
        \begin{enumerate}[label=\textbf{問\arabic*.},leftmargin=3.8em,format=\bookitemlink{practice-advanced}{solution}{problem}]
          % A1: 階乗で重みをそろえた総和
          \item まず総和を消そうとすると,添字によって階乗も係数も変わる.そこで,総和の各項に合わせて
                \[
                  b_n=\frac{a_n}{n!},\qquad
                  T_n=\sum_{k=1}^{n}b_k
                \]
                とおく.両辺を$(n+1)!$で割れば,
                \[
                  b_{n+1}=4b_n-T_n,\qquad b_1=1.
                \]
                これより$b_2=3$である.次の式では
                \[
                  b_{n+2}=4b_{n+1}-T_{n+1}
                          =3b_{n+1}-T_n
                \]
                となるので,\, $T_n=4b_n-b_{n+1}$を代入して
                \[
                  b_{n+2}-4b_{n+1}+4b_n=0
                \]
                を得る.特性方程式は$(r-2)^2=0$であり,同じ等比数列を2本作っても情報は増えない.ここでは
                \[
                  c_n=\frac{b_n}{2^{n-1}}
                \]
                とおくと,
                \[
                  c_{n+2}-2c_{n+1}+c_n=0.
                \]
                したがって$c_{n+1}-c_n$は一定である.初めの2項は$c_1=1$,\, $c_2=3/2$だから,
                \[
                  c_n=1+\frac{n-1}{2}=\frac{n+1}{2}.
                \]
                よって$b_n=(n+1)2^{n-2}$となり,\, $a_n=n!b_n$から一般項が得られる.階乗による正規化を最初に行うと,その後は本編の総和型と重解型の手順で解ける.
          % A2: 添字を含む分母と不動点
          \item 分子と分母の係数を個別に追うより,差を計算しよう.
                \[
                  a_{n+1}-1
                  =\frac{a_n-1}{2+n(a_n-1)}.
                \]
                $a_1>1$であり,\, $a_n>1$ならば右辺も正である.したがって帰納的に$a_n>1$が成り立ち,元の分母も常に正となる.
                
                ここで
                \[
                  b_n=\frac{1}{a_n-1}
                \]
                とおくと,
                \[
                  b_{n+1}=2b_n+n,\qquad b_1=1.
                \]
                逆数を取ると分数型は消えるが,まだ等比型ではない.右辺の$n$を消すため,
                \[
                  c_n=b_n+n+1
                \]
                とおけば,
                \begin{align*}
                  c_{n+1}
                    &=2b_n+n+(n+1)+1\\
                    &=2(b_n+n+1)=2c_n.
                \end{align*}
                $c_1=3$だから$c_n=3\cdot2^{n-1}$であり,
                \[
                  b_n=3\cdot2^{n-1}-n-1.
                \]
                逆数を戻して一般項を得る.なお,\, $b_1=1$で,\, $b_{n+1}=2b_n+n>0$なので,得られた式の分母も$0$にならない.
                
                この問題では,不動点への移動,逆数化,一次式の移動を別々の目的で使っている.一度の変換で解ける形にならなくても,次に消すべき項を見極めればよい.

                \noindent 【別解】\\
                $b_n$の式を$2^n$で割ると,
                \[
                 \frac{b_{n+1}}{2^n}
                  -\frac{b_n}{2^{n-1}}=\frac{n}{2^n}.
                \]
                したがって
                \[
                 \frac{b_n}{2^{n-1}}
                 =1+\sum_{k=1}^{n-1}\frac{k}{2^k}.
                \]
                ここで
                \[
                 \frac{k}{2^k}
                 =\frac{k+1}{2^{k-1}}-\frac{k+2}{2^k}
                \]
                だから総和は相殺し,
                \[
                 \sum_{k=1}^{n-1}\frac{k}{2^k}
                  =2-\frac{n+1}{2^{n-1}}.
                \]
                これから同じ$b_n$を得る.一次式を加える解法は,この相殺を補助数列の形で先に組み込んだものといえる.
          % A3: 三角形状の重みをもつ積
          \item 各項は正なので,
                \[
                  x_n=\log_2a_n
                \]
                とおける.対数を取ると積が和になり,
                \[
                  x_{n+1}
                   =n+\sum_{k=1}^{n}(n-k+1)x_k,
                  \qquad x_1=0
                \]
                となる.この式から$x_2=1$である.
                
                重みのある総和を一度に計算する必要はない.同じ形の式を1つずらして引くと,\, $n\ge2$について
                \[
                  x_{n+1}-x_n
                  =1+\sum_{k=1}^{n}x_k.
                \]
                $n=1$でも$x_2-x_1=1+x_1$が成り立つので,この関係はすべての$n\ge1$で使える.さらに隣接する式の差を取ると,
                \[
                  x_{n+2}-2x_{n+1}+x_n=x_{n+1}.
                \]
                したがって,
                \[
                  x_{n+2}-3x_{n+1}+x_n=0.
                \]
                特性方程式の2根を
                \[
                  \alpha=\frac{3+\sqrt5}{2},\qquad
                  \beta=\frac{3-\sqrt5}{2}
                \]
                とすると,\, $\alpha+\beta=3$,\, $\alpha\beta=1$である.よって
                \[
                  x_{n+2}-\beta x_{n+1}
                   =\alpha(x_{n+1}-\beta x_n).
                \]
                初期値$x_1=0$,\, $x_2=1$に合う形は
                \[
                  x_n
                   =\frac{\alpha^{n-1}-\beta^{n-1}}{\alpha-\beta}
                   =\frac{\alpha^{n-1}-\beta^{n-1}}{\sqrt5}.
                \]
                実際,この式は上の3項間関係と2つの初期値を満たす.その関係から次項は一意に決まるので,これが$x_n$である.最後に$a_n=2^{x_n}$と戻せばよい.
                
                指数は$0,1,3,8,21,\ldots$となる.巨大な積を計算する難しさは,重み付き総和の差を2回取ることで消えている.
          % A4: 非線形3項間式から階比を作る
          \item 元の式は項の積を含むが,次数はすべて$2$である.各項は正なので,
                \[
                  r_n=\frac{a_{n+1}}{a_n}
                \]
                とおき,元の式を$a_na_{n+1}$で割ると,
                \[
                  (n+r_n)r_{n+1}=(n+1)r_n.
                \]
                したがって,
                \[
                  r_{n+1}=\frac{(n+1)r_n}{n+r_n},
                  \qquad r_1=2.
                \]
                ここで逆数$c_n=1/r_n$を取れば,
                \[
                  c_{n+1}=\frac{nc_n+1}{n+1}.
                \]
                各$n$で$c_n=1$が変わらないことに注目して,
                \[
                  d_n=c_n-1
                \]
                とおくと,
                \[
                  d_{n+1}=\frac{n}{n+1}d_n,\qquad d_1=-\frac12.
                \]
                よって積の相殺から
                \[
                  d_n=-\frac1{2n},\qquad
                  c_n=\frac{2n-1}{2n},\qquad
                  r_n=\frac{2n}{2n-1}.
                \]
                初項$a_1=1$に戻って,
                \[
                  a_n=\prod_{k=1}^{n-1}\frac{2k}{2k-1}.
                \]
                $m=n-1$とおくと分子は$2^m m!$である.分母の奇数の積については,
                \[
                  (2m)!
                   =(1\cdot3\cdots(2m-1))\,2^m m!
                \]
                なので,
                \[
                  a_n=\frac{4^m(m!)^2}{(2m)!}.
                \]
                $m=0$では空積を$1$とし,\, $a_1=1$も含められる.分数型を解く前に,階比を導入することが決め手である.

                \noindent 【別解】\\
                逆数の式で,平行移動せず直接両辺に$n+1$を掛けてもよい.
                \[
                  (n+1)c_{n+1}-nc_n=1.
                \]
                したがって$e_n=nc_n$は公差$1$の等差数列であり,\, $e_1=1/2$から
                \[
                  e_n=n-\frac12,\qquad c_n=1-\frac1{2n}
                \]
                となる.不動点を引く方法と,添字を掛けて階差を一定にする方法の両方が使える.
          % A5: 和と差で異なる構造が現れる連立
          \item 2式で共通する部分は,引けば消える.符号だけが違う部分は,足せば消える.そこで
                \[
                  s_n=a_n+b_n,\qquad d_n=a_n-b_n
                \]
                とおくと,
                \[
                  s_{n+1}=2(n+1)s_n,\qquad s_1=2,
                \]
                \[
                  d_{n+1}
                   =\frac{3d_n}{3+(n+1)d_n},\qquad d_1=1.
                \]
                $d_n>0$なら$d_{n+1}>0$であるから,帰納的に差は常に正となり,元の分母は$0$にならない.
                
                和の式は,階乗で割ると
                \[
                  \frac{s_{n+1}}{(n+1)!}=2\frac{s_n}{n!}
                \]
                となる.よって
                \[
                  s_n=2^n n!.
                \]
                一方,差の式の逆数を取って$c_n=1/d_n$とおけば,
                \[
                  c_{n+1}=c_n+\frac{n+1}{3},\qquad c_1=1.
                \]
                したがって,
                \begin{align*}
                 c_n
                  &=1+\frac13\sum_{k=1}^{n-1}(k+1)\\
                  &=1+\frac{(n-1)(n+2)}6
                   =\frac{n^2+n+4}{6}.
                \end{align*}
                これより$d_n=6/(n^2+n+4)$である.最後に
                \[
                  a_n=\frac{s_n+d_n}{2},\qquad
                  b_n=\frac{s_n-d_n}{2}
                \]
                と戻して一般項を得る.
                
                連立式全体を一つの方法で解こうとする必要はない.和と差を作った後,それぞれに適した解法を選ぶ.
          % A6: 積の3項間関係と二重の階差
          \item 正値の条件を使って$x_n=\log_2a_n$とおくと,
                \[
                  x_{n+2}-2x_{n+1}+x_n=2n+1,
                  \qquad x_1=1,\quad x_2=4.
                \]
                この式は$x_n$自体の階差ではなく,階差の階差を表している.そこで
                \[
                  d_n=x_{n+1}-x_n
                \]
                とおけば,
                \[
                  d_{n+1}-d_n=2n+1,\qquad d_1=3.
                \]
                $2n+1=(n+1)^2-n^2$なので,
                \[
                  d_n=3+(n^2-1)=n^2+2.
                \]
                さらにもう一度階差を加えて,
                \begin{align*}
                  x_n
                   &=1+\sum_{k=1}^{n-1}(k^2+2)\\
                   &=1+\frac{n(n-1)(2n-1)}6+2(n-1)\\
                   &=\frac{n(n-1)(2n-1)}6+2n-1.
                \end{align*}
                最後に$a_n=2^{x_n}$と戻せばよい.
                
                本編の隣接3項間漸化式と似た形でも,定数項が$n$に依存し,しかも左辺が二重の階差そのものになっている.この構造を見れば,特性方程式の計算に進む前に短い解法が見つかる.

                \noindent 【別解】\\
                $r_n=a_{n+1}/a_n$とおき,元の式を$a_na_{n+1}$で割ると,
                \[
                  r_{n+1}=2^{2n+1}r_n,\qquad r_1=8.
                \]
                これは階比型なので,
                \begin{align*}
                 r_n
                  &=2^{\,3+\sum_{k=1}^{n-1}(2k+1)}\\
                  &=2^{\,n^2+2}.
                \end{align*}
                もう一度比を掛け合わせて,
                \[
                 a_n
                  =2\prod_{k=1}^{n-1}2^{k^2+2}
                  =2^{\,1+\sum_{k=1}^{n-1}(k^2+2)}
                \]
                を得る.対数で差を取る方法と,元の数列で比を取る方法は,同じ構造を異なる形で使っている.
          % A7: 添字とともに動く基準点
          \item $a_n=n$なら$a_{n+1}=n+1$となるので,基準になる値も$n$に応じて変わっている.そこで
                \[
                  c_n=\frac{a_n}{n}
                \]
                とおくと,
                \[
                  c_{n+1}
                   =\frac{(n+1)c_n+1}{c_n+n+1},
                  \qquad c_1=3.
                \]
                $c_n>1$なら分母は正であり,
                \[
                  c_{n+1}-1=\frac{n(c_n-1)}{c_n+n+1}>0.
                \]
                したがって全項$c_n>1$が成り立ち,元の分母も正となる.
                
                この式では$c_n=1$と$c_n=-1$がともに動かない.両方の距離を計算すると,
                \[
                  c_{n+1}+1
                   =\frac{(n+2)(c_n+1)}{c_n+n+1}.
                \]
                同じ分母が消えるように,
                \[
                  b_n=\frac{c_n-1}{c_n+1}
                \]
                とおけば,
                \[
                  b_{n+1}=\frac{n}{n+2}b_n,\qquad b_1=\frac12.
                \]
                よって,
                \[
                  b_n
                   =\frac12\prod_{k=1}^{n-1}\frac{k}{k+2}
                   =\frac{1}{n(n+1)}.
                \]
                $c_n=(1+b_n)/(1-b_n)$だから,
                \[
                  c_n=\frac{n(n+1)+1}{n(n+1)-1}.
                \]
                最後に$a_n=nc_n$と戻す.ここでは正規化によって分数型を見つけ,さらにその分数型を階比型へ変えている.

                \noindent 【別解】\\
                $c_n$の式から$u_n=1/(c_n-1)$とおいてもよい.
                \[
                 u_{n+1}=\frac{n+2}{n}u_n+\frac1n.
                \]
                定数$1/2$を加えると,
                \[
                 u_{n+1}+\frac12
                   =\frac{n+2}{n}\left(u_n+\frac12\right).
                \]
                $u_1=1/2$から,\, $u_n+1/2=n(n+1)/2$を得る.二基準点の比を取る方法と,一基準点との差の逆数を取ってから定数を加える方法が対応している.
          % A8: 係数の和が零になる3項間関係
          \item 係数は$n$に依存するため,定数係数の特性方程式型ではない.一方,
                \[
                  2n^2+10n+11
                   =(n+2)(n+4)+(n+1)(n+3)
                \]
                なので,元の式は
                \begin{align*}
                 &(n+2)(n+4)(a_{n+2}-a_{n+1})\\
                 &\qquad=(n+1)(n+3)(a_{n+1}-a_n)
                \end{align*}
                とまとめられる.ここで$d_n=a_{n+1}-a_n$とおけば,
                \[
                  d_{n+1}
                   =\frac{(n+1)(n+3)}{(n+2)(n+4)}d_n,
                  \qquad d_1=1.
                \]
                2つの比に分けて掛け合わせると,
                \begin{align*}
                 d_n
                  &=\prod_{k=1}^{n-1}
                      \frac{k+1}{k+2}\frac{k+3}{k+4}\\
                  &=\frac2{n+1}\frac4{n+3}
                   =\frac8{(n+1)(n+3)}.
                \end{align*}
                次に元の数列を求めるには,この階差を足す.
                \[
                 d_k=4\left(\frac1{k+1}-\frac1{k+3}\right)
                \]
                だから,\, $n\ge2$について
                \begin{align*}
                 a_n
                  &=4\sum_{k=1}^{n-1}
                       \left(\frac1{k+1}-\frac1{k+3}\right)\\
                  &=4\left(\frac12+\frac13
                          -\frac1{n+1}-\frac1{n+2}\right).
                \end{align*}
                これを整理すれば一般項が得られる.得られた式は$n=1$でも$a_1=0$を満たす.階差を作って終わるのでなく,差の階比を解いた後,総和で戻るところまでが一連の手順である.

                \noindent 【別解】\\
                階差の式に両辺の係数を掛ける形のまま,
                \[
                 (n+2)(n+4)d_{n+1}
                   =(n+1)(n+3)d_n
                \]
                を眺めると,\, $e_n=(n+1)(n+3)d_n$が定数数列である.初項$e_1=8$から直ちに$d_n=8/\{(n+1)(n+3)\}$を得る.相殺する積を全部書く代わりに,相殺後に残る因子を補助数列へ組み込める.
          % A9: 掛け算で結び付いた二つの数列
          \item 各項は正なので,
                \[
                  x_n=\log_2a_n,\qquad y_n=\log_2b_n
                \]
                とおくと,
                \[
                 \begin{cases}
                  x_{n+1}=2x_n+y_n+2n+1,\\
                  y_{n+1}=x_n+2y_n-2n-1
                 \end{cases}
                \]
                となる.ここからは本編の連立型と同じように,
                \[
                  s_n=x_n+y_n,\qquad d_n=x_n-y_n
                \]
                を使う.足し算では$n$の項が消えて,
                \[
                  s_{n+1}=3s_n,\qquad s_1=2,
                \]
                引き算では
                \[
                  d_{n+1}=d_n+4n+2,\qquad d_1=0
                \]
                となる.よって$s_n=2\cdot3^{n-1}$である.また
                \[
                 4n+2=2\{(n+1)^2-n^2\}
                \]
                なので,
                \[
                 d_n=2(n^2-1).
                \]
                したがって,
                \[
                 x_n=\frac{s_n+d_n}{2}=3^{n-1}+n^2-1,
                \]
                \[
                 y_n=\frac{s_n-d_n}{2}=3^{n-1}-n^2+1.
                \]
                指数に戻して答えを得る.対数型と連立型は別々の分類に見えるが,同じ問題で順に使うことができる.

                \noindent 【別解】\\
                対数を取らずに,積と商を作っても分離できる.
                \[
                  P_n=a_nb_n,\qquad R_n=\frac{a_n}{b_n}
                \]
                とおくと,
                \[
                  P_{n+1}=P_n^3,\quad P_1=4,
                  \qquad R_{n+1}=2^{4n+2}R_n,\quad R_1=1.
                \]
                よって,
                \[
                  P_n=2^{\,2\cdot3^{n-1}},\qquad
                  R_n=2^{\,2(n^2-1)}.
                \]
                正値を使って$a_n=\sqrt{P_nR_n}$,\, $b_n=\sqrt{P_n/R_n}$と戻せば同じ答えになる.対数の和差は,元の数列の積と商に対応する.
          % A10: 二つの平方が作る分数式
          \item 分子だけを展開すると構造が見えにくくなる.分母を$D_n$と書けば,
                \[
                 a_{n+1}-1=\frac{2^{n+1}(a_n-1)^2}{D_n},
                 \qquad
                 a_{n+1}+1=\frac{2(a_n+1)^2}{D_n}.
                \]
                共通する分母を消すため,
                \[
                  b_n=\frac{a_n-1}{a_n+1}
                \]
                とおくと,
                \[
                  b_{n+1}=2^n b_n^2,\qquad b_1=\frac18.
                \]
                ここで,次項が定義できることも帰納的に確かめておこう.
                \[
                  a_n>1,\qquad 0<b_n\le2^{-n-2}
                \]
                が$n=1$で成り立つ.この評価がある$n$で成り立つと仮定すると,まず
                \[
                  D_n=(a_n+1)^2(1-2^nb_n^2)>0
                \]
                なので,\, $a_{n+1}$が定義できる.また$a_{n+1}-1$の式は正であり,\, $a_{n+1}>1$である.その後,比の式を使えば
                \[
                  0<b_{n+1}\le2^{-n-4}<2^{-(n+1)-2}
                \]
                を得る.したがって,元の数列はすべての$n$で定義でき,逆数や対数の操作も有効である.
                
                次に$x_n=-\log_2b_n$とおけば,
                \[
                  x_{n+1}=2x_n-n,\qquad x_1=3.
                \]
                $n$の一次式を消すため$y_n=x_n-n-1$とおくと,
                \[
                  y_{n+1}=2y_n,\qquad y_1=1.
                \]
                よって$x_n=n+1+2^{n-1}=E_n$であり,
                \[
                  b_n=2^{-E_n}.
                \]
                $a_n=(1+b_n)/(1-b_n)$と戻して答えを得る.分数の変形だけでも,対数だけでも解き切れない.平方の構造を作ってから対数を使う順序が大切である.
        \end{enumerate}

